% Source: https://artofproblemsolving.com/community/c3h1346754_motion_problem
% Stable identity: aops-topic-1346754-q01
Downstream a vessel moves at (still-water speed) $+$ (current), upstream at (still-water speed) $-$ (current).

\textbf{The current, from the yacht.} The yacht's downstream speed is $\dfrac{90}{6}=15$ mph and its upstream speed is $\dfrac{90}{10}=9$ mph. Their difference is twice the current, so the current is $\dfrac{15-9}{2}=3$ mph.

\textbf{The boat.} Its downstream speed is $\dfrac{90}{5}=18$ mph, so its still-water speed is $18-3=15$ mph and its upstream speed is $15-3=12$ mph. The trip upstream takes
\[
\dfrac{90}{12}=\boxed{7.5\text{ hours}}.
\]
